% ------------------------------------------------------------------
\begin{frame}{今日のゴール / Today's Goal}
  \begin{block}{目標}
    ws:: API 全 43 関数をカテゴリ別に理解し、各関数の使い方を把握する
  \end{block}

  \vspace{0.5em}

  \begin{itemize}
    \item ws:: API 全 43 関数をカテゴリ別に理解
    \item 各 API の使い方をコード例で確認
    \item 環境・距離センサ API で全センサに直接アクセス
  \end{itemize}
\end{frame}

% ==================================================================
%   ws:: API Category Pages
% ==================================================================

% ------------------------------------------------------------------
\begin{frame}{ws:: API 概要 / API Overview}
  \centering
  {\small
  \begin{tabular}{llr}
    \hline
    \textbf{カテゴリ} & \textbf{概要} & \textbf{関数数} \\
    \hline
    Motor Control   & モータ制御・Arm/Disarm    & 7 \\
    RC Input        & コントローラスティック入力 & 4 \\
    RC Buttons      & ボタン・飛行モード         & 5 \\
    IMU Sensor      & ジャイロ・加速度センサ     & 6 \\
    Env / Distance  & 気圧・磁気・ToF・光学フロー & 10 \\
    Estimation      & ESKF 姿勢・高度推定        & 4 \\
    LED             & LED 色制御                 & 3 \\
    Utility         & 時刻・電圧・通信・出力     & 4 \\
    \hline
    \multicolumn{2}{l}{\textbf{合計}} & \textbf{43} \\
    \hline
  \end{tabular}
  }

  \vspace{0.5em}
  {\footnotesize
    \texttt{\#include "workshop\_api.hpp"} で全関数を使用可能。
    全関数は \texttt{ws::} 名前空間に属する。\\
    より深いアクセスには \textbf{StampFlyState} を使用（後半で解説）。
  }
\end{frame}

% ------------------------------------------------------------------
\begin{frame}[fragile]{Motor Control API --- Duty \& Mixer}
  \vspace{0.3em}
  {\small
  \begin{tabular}{lll}
    \hline
    \textbf{関数} & \textbf{引数} & \textbf{説明} \\
    \hline
    \texttt{motor\_set\_duty(id, duty)} & id: 1--4, duty: 0--1 & 個別モータ duty 設定 \\
    \texttt{motor\_set\_all(duty)}      & duty: 0--1           & 全モータ同一 duty \\
    \texttt{motor\_stop\_all()}         & ---                   & 全モータ即時停止 \\
    \texttt{motor\_mixer(t, r, p, y)}   & 各 float             & スラスト+姿勢ミキシング \\
    \hline
  \end{tabular}
  }

  \vspace{0.3em}

  \begin{lstlisting}[style=sfcpp, basicstyle=\ttfamily\scriptsize]
// Individual motor control
ws::motor_set_duty(1, 0.3f);  // Motor 1 (FR) = 30%
ws::motor_stop_all();          // Emergency stop

// Mixer: thrust + attitude corrections
ws::motor_mixer(0.5f, 0.0f, 0.0f, 0.0f);  // Hover
  \end{lstlisting}
\end{frame}

% ------------------------------------------------------------------
\begin{frame}[fragile]{Motor Control API --- Arm / Disarm}
  \vspace{0.3em}
  {\small
  \begin{tabular}{lll}
    \hline
    \textbf{関数} & \textbf{戻り値} & \textbf{説明} \\
    \hline
    \texttt{arm()}       & void & モータ出力を有効化（Arm） \\
    \texttt{disarm()}    & void & モータ出力を無効化（Disarm） \\
    \texttt{is\_armed()} & bool & Arm 状態を返す（\texttt{true} = Armed） \\
    \hline
  \end{tabular}
  }

  \vspace{0.3em}

  \begin{lstlisting}[style=sfcpp, basicstyle=\ttfamily\scriptsize]
if (ws::rc_throttle() < 0.05f) ws::arm();

if (ws::is_armed()) {
    ws::motor_mixer(ws::rc_throttle(),
        ws::rc_roll(), ws::rc_pitch(), ws::rc_yaw());
} else {
    ws::motor_stop_all();
}
  \end{lstlisting}

  \vspace{0.2em}
  {\footnotesize \warn{注意:} \api{arm()} 前はモータコマンドを送っても回転しない}
\end{frame}

% ------------------------------------------------------------------
\begin{frame}[fragile]{RC Input API}
  \vspace{0.3em}
  {\small
  \begin{tabular}{lll}
    \hline
    \textbf{関数} & \textbf{範囲} & \textbf{説明} \\
    \hline
    \texttt{rc\_throttle()} & $[0.0, 1.0]$   & スロットル（上がプラス） \\
    \texttt{rc\_roll()}     & $[-1.0, 1.0]$  & ロール（右がプラス） \\
    \texttt{rc\_pitch()}    & $[-1.0, 1.0]$  & ピッチ（前がプラス） \\
    \texttt{rc\_yaw()}      & $[-1.0, 1.0]$  & ヨー（右回転がプラス） \\
    \hline
  \end{tabular}
  }

  \vspace{0.3em}

  \begin{lstlisting}[style=sfcpp, basicstyle=\ttfamily\scriptsize]
void loop_400Hz(float dt) {
    float thr = ws::rc_throttle();  // 0.0 - 1.0
    float rol = ws::rc_roll();      // -1.0 - 1.0
    float pit = ws::rc_pitch();
    float yaw = ws::rc_yaw();
    ws::motor_mixer(thr, rol, pit, yaw);  // ACRO mode
}
  \end{lstlisting}
\end{frame}

% ------------------------------------------------------------------
\begin{frame}[fragile]{RC Buttons / Mode API}
  \vspace{0.3em}
  {\small
  \begin{tabular}{lll}
    \hline
    \textbf{関数} & \textbf{戻り値} & \textbf{説明} \\
    \hline
    \texttt{rc\_throttle\_yaw\_button()}  & bool & 左スティック押し込みボタン \\
    \texttt{rc\_roll\_pitch\_button()}    & bool & 右スティック押し込みボタン \\
    \texttt{rc\_stabilize\_acro\_mode()}  & bool & ACRO モード判定（\texttt{true} = ACRO） \\
    \texttt{rc\_alt\_mode()}             & bool & 高度保持モード判定 \\
    \texttt{rc\_pos\_mode()}             & bool & 位置保持モード判定 \\
    \hline
  \end{tabular}
  }

  \vspace{0.3em}

  \begin{lstlisting}[style=sfcpp, basicstyle=\ttfamily\scriptsize]
void loop_400Hz(float dt) {
    if (ws::rc_stabilize_acro_mode()) {
        acro_control(dt);       // ACRO: rate control
    } else {
        stabilize_control(dt);  // Stabilize: angle control
    }
}
  \end{lstlisting}
\end{frame}

% ------------------------------------------------------------------
\begin{frame}[fragile]{IMU Sensor API}
  \vspace{0.3em}
  {\small
  \begin{tabular}{lll}
    \hline
    \textbf{関数} & \textbf{単位} & \textbf{説明} \\
    \hline
    \texttt{gyro\_x/y/z()}  & rad/s   & R/P/Y 角速度（BMI270） \\
    \texttt{accel\_x/y/z()} & m/s$^2$ & X/Y/Z 加速度（NED座標系） \\
    \hline
  \end{tabular}
  }

  \vspace{0.3em}

  \begin{lstlisting}[style=sfcpp, basicstyle=\ttfamily\scriptsize]
void loop_400Hz(float dt) {
    float gx = ws::gyro_x();    // [rad/s] roll rate
    float gy = ws::gyro_y();    // [rad/s] pitch rate
    float gz = ws::gyro_z();    // [rad/s] yaw rate
    float az = ws::accel_z();   // [m/s^2] ~9.81 at rest

    ws::print(">gyro_x:%.3f", gx);
    ws::print(">accel_z:%.2f", az);
}
  \end{lstlisting}

  \vspace{0.2em}
  {\footnotesize 静止時: \texttt{accel\_z()} $\approx$ 9.81（NED 座標系、下向き正）}
\end{frame}

% ------------------------------------------------------------------
\begin{frame}[fragile]{Estimation API}
  \vspace{0.3em}
  {\small
  \begin{tabular}{lll}
    \hline
    \textbf{関数} & \textbf{単位} & \textbf{説明} \\
    \hline
    \texttt{estimated\_roll()}     & rad & ESKF ロール推定角 \\
    \texttt{estimated\_pitch()}    & rad & ESKF ピッチ推定角 \\
    \texttt{estimated\_yaw()}      & rad & ESKF ヨー推定角 \\
    \texttt{estimated\_altitude()} & m   & ESKF 推定高度（正 = 上） \\
    \hline
  \end{tabular}
  }

  \vspace{0.3em}

  \begin{lstlisting}[style=sfcpp, basicstyle=\ttfamily\scriptsize]
void loop_400Hz(float dt) {
    float roll  = ws::estimated_roll();     // [rad]
    float pitch = ws::estimated_pitch();
    float alt   = ws::estimated_altitude(); // [m]
    // P control: level the drone
    ws::motor_mixer(ws::rc_throttle(),
        (0.0f-roll)*0.5f, (0.0f-pitch)*0.5f, ws::rc_yaw());
}
  \end{lstlisting}

  \vspace{0.2em}
  {\footnotesize ESKF はジャイロ+加速度+気圧+ToF+磁気+光学フローを統合（L09 参照）}
\end{frame}

% ------------------------------------------------------------------
\begin{frame}[fragile]{環境センサ API --- 気圧・磁気}
  \vspace{0.3em}
  {\small
  \begin{tabular}{lll}
    \hline
    \textbf{関数} & \textbf{単位} & \textbf{説明} \\
    \hline
    \texttt{baro\_altitude()} & m   & 気圧高度（BMP280） \\
    \texttt{baro\_pressure()} & Pa  & 気圧値（BMP280） \\
    \texttt{mag\_x()}         & $\mu$T & 磁気 X 成分（BMM150） \\
    \texttt{mag\_y()}         & $\mu$T & 磁気 Y 成分（BMM150） \\
    \texttt{mag\_z()}         & $\mu$T & 磁気 Z 成分（BMM150） \\
    \hline
  \end{tabular}
  }

  \vspace{0.3em}

  \begin{lstlisting}[style=sfcpp, basicstyle=\ttfamily\scriptsize]
float alt = ws::baro_altitude();  // [m]
float p   = ws::baro_pressure();  // [Pa]
float mx  = ws::mag_x();          // [uT]
ws::print(">baro_alt:%.2f", alt);
  \end{lstlisting}
\end{frame}

% ------------------------------------------------------------------
\begin{frame}[fragile]{環境センサ API --- ToF・光学フロー}
  \vspace{0.3em}
  {\small
  \begin{tabular}{lll}
    \hline
    \textbf{関数} & \textbf{単位} & \textbf{説明} \\
    \hline
    \texttt{tof\_bottom()}   & m     & 下方 ToF 距離（VL53L3CX, 0--2\,m） \\
    \texttt{tof\_front()}    & m     & 前方 ToF 距離（$-1$ = 未接続） \\
    \texttt{flow\_vx()}      & m/s   & 光学フロー速度 X（PMW3901） \\
    \texttt{flow\_vy()}      & m/s   & 光学フロー速度 Y（PMW3901） \\
    \texttt{flow\_quality()} & 0--255 & 光学フロー品質（高い = 良好） \\
    \hline
  \end{tabular}
  }

  \vspace{0.3em}

  \begin{lstlisting}[style=sfcpp, basicstyle=\ttfamily\scriptsize]
float dist = ws::tof_bottom();    // [m] ground distance
float vx   = ws::flow_vx();       // [m/s]
uint8_t sq = ws::flow_quality();   // 0-255
ws::print(">tof_bottom:%.3f", dist);
  \end{lstlisting}
\end{frame}

% ------------------------------------------------------------------
\begin{frame}[fragile]{LED Control API}
  \vspace{0.3em}
  {\small
  \begin{tabular}{lll}
    \hline
    \textbf{関数} & \textbf{引数/戻り値} & \textbf{説明} \\
    \hline
    \texttt{led\_color(r, g, b)}      & r,g,b: 0--255 & LED 色を RGB 設定 \\
    \texttt{disable\_led\_task()}     & void           & システム LED タスク停止 \\
    \texttt{enable\_led\_task()}      & void           & システム LED タスク再開 \\
    \texttt{is\_led\_task\_disabled()} & bool          & LED タスク停止中か確認 \\
    \hline
  \end{tabular}
  }

  \vspace{0.3em}

  \begin{lstlisting}[style=sfcpp, basicstyle=\ttfamily\scriptsize]
void setup() {
    ws::disable_led_task();  // Take over LED control
}
void loop_400Hz(float dt) {
    float v = ws::battery_voltage();
    if (v > 3.8f)      ws::led_color(0, 255, 0);   // Green
    else if (v > 3.5f) ws::led_color(255, 165, 0); // Orange
    else                ws::led_color(255, 0, 0);   // Red
}
  \end{lstlisting}

  \vspace{0.2em}
  {\footnotesize \warn{注意:} \api{disable\_led\_task()} を先に呼ぶこと（LED 上書き防止）}
\end{frame}

% ------------------------------------------------------------------
\begin{frame}[fragile]{Utility API}
  \vspace{0.3em}
  {\small
  \begin{tabular}{lll}
    \hline
    \textbf{関数} & \textbf{戻り値} & \textbf{説明} \\
    \hline
    \texttt{millis()}           & uint32\_t (ms) & 起動からの経過時間 \\
    \texttt{battery\_voltage()} & float (V)      & バッテリー電圧 \\
    \texttt{print(fmt, ...)}    & void           & printf 形式でデバッグ出力 \\
    \texttt{set\_channel(ch)}   & void           & WiFi チャネル設定（1, 6, 11） \\
    \hline
  \end{tabular}
  }

  \vspace{0.3em}

  \begin{lstlisting}[style=sfcpp, basicstyle=\ttfamily\scriptsize]
void setup() { ws::set_channel(6); }
void loop_400Hz(float dt) {
    static uint32_t last = 0;
    uint32_t now = ws::millis();
    if (now - last >= 20) {  // 50 Hz decimation
        last = now;
        ws::print(">battery:%.2f", ws::battery_voltage());
    }
}
  \end{lstlisting}
\end{frame}

% ==================================================================
%   Advanced: StampFlyState
% ==================================================================

% ------------------------------------------------------------------
\begin{frame}{ハードウェアセンサ仕様 / Hardware Sensors}
  \centering

  {\small
  \begin{tabular}{llll}
    \hline
    \textbf{センサ} & \textbf{型番} & \textbf{サンプルレート} & \textbf{測定量} \\
    \hline
    IMU        & BMI270  & 400\,Hz  & 加速度 + ジャイロ \\
    気圧       & BMP280  & 50\,Hz   & 気圧 → 高度 \\
    磁気       & BMM150  & 10--30\,Hz & 磁気ベクトル \\
    ToF（下方） & VL53L3CX & 30\,Hz  & 対地距離（0--2\,m） \\
    ToF（前方） & VL53L3CX & 30\,Hz  & 前方距離（0--2\,m） \\
    光学フロー  & PMW3901 & 100\,Hz  & 対地速度 \\
    \hline
  \end{tabular}
  }

  \vspace{1em}

  {\footnotesize
    ws:: API で全センサに直接アクセス可能（\api{baro\_altitude()}, \api{tof\_bottom()} 等）。\\
    より低レベルなアクセスには \textbf{StampFlyState} を使用。\\
    ESKF は全センサを統合して姿勢・位置・速度を推定（Lesson 9 参照）。
  }
\end{frame}

% ------------------------------------------------------------------
\begin{frame}[fragile]{Teleplot}
  {\small Teleplot 形式 (\texttt{>name:value}) でシリアル出力すると、\\
  VS Code 拡張 Teleplot がリアルタイムグラフを描画する。}

  \vspace{0.3em}

  \begin{lstlisting}[style=sfcpp, basicstyle=\ttfamily\scriptsize]
void loop_400Hz(float dt) {
    static uint32_t tick = 0; tick++;
    if (tick % 8 != 0) return;  // 50 Hz decimation

    // Teleplot output (>name:value format)
    ws::print(">baro_alt:%.2f", ws::baro_altitude());
    ws::print(">tof_bottom:%.3f", ws::tof_bottom());
    ws::print(">eskf_alt:%.2f", ws::estimated_altitude());
}
  \end{lstlisting}

  \vspace{0.2em}
  {\footnotesize
    \textbf{セットアップ:} VSCode 拡張 \texttt{alexnesnes.teleplot} をインストール
    → Teleplot パネルでシリアルポートを選択 → 自動グラフ化
  }
\end{frame}

% ------------------------------------------------------------------
\begin{frame}[fragile]{実習: 全センサ Teleplot 出力 / Hands-on}
  \begin{lstlisting}[style=sfcpp, basicstyle=\ttfamily\scriptsize]
#include "workshop_api.hpp"
void setup() { ws::print("Lesson 10: Sensor API"); }
void loop_400Hz(float dt) {
    static uint32_t tick = 0; tick++;
    if (tick % 8 != 0) return;  // 50 Hz decimation
    // Environmental sensors (ws:: API)
    ws::print(">baro_alt:%.2f", ws::baro_altitude());
    ws::print(">mag_x:%.1f", ws::mag_x());
    ws::print(">tof_bottom:%.3f", ws::tof_bottom());
    ws::print(">flow_vx:%.3f", ws::flow_vx());
    // ESKF estimation
    ws::print(">eskf_alt:%.2f", ws::estimated_altitude());
    ws::print(">eskf_roll:%.1f", ws::estimated_roll()*57.3f);
}
  \end{lstlisting}

  \vspace{0.2em}
  {\footnotesize
    Teleplot パネルでシリアルポートを選択し確認（前ページ参照）
  }
\end{frame}

% ------------------------------------------------------------------
\begin{frame}{チェックポイント / Checkpoint}
  \begin{alertblock}{確認事項}
    \begin{itemize}
      \item[$\square$] ws:: API 全 43 関数のカテゴリと役割を把握した
      \item[$\square$] Motor / RC / IMU / Sensor / Estimation の使い方を理解した
      \item[$\square$] 気圧・磁気・ToF・光学フロー API で値を取得できた
      \item[$\square$] Teleplot で複数センサのグラフを同時表示した
    \end{itemize}
  \end{alertblock}

  \vspace{0.5em}

  \begin{block}{次のレッスン}
    Lesson 11: 独自ファームウェア開発
  \end{block}
\end{frame}
